% Standalone editable source for Molarium Appendix B.
% Compile with: latexmk -pdf Molarium_Appendix_B_Workflows_for_Humans_and_Agents.tex
%
% Editorial contract for future agents:
% 1. Organize by the scientific task/question, not by software module.
% 2. Pair each human workflow with an explicit agent skill contract.
% 3. State what changes the authoritative molecular state and what remains a candidate.
% 4. Preserve implementation status, evidence, failure behavior, and scientific boundaries.
% 5. Do not silently turn repository-only or proposed functionality into a current workflow.

\documentclass[10pt,letterpaper]{article}

\usepackage[T1]{fontenc}
\usepackage{lmodern}
\usepackage[letterpaper,left=56.2pt,right=56.2pt,top=51.8pt,bottom=41.8pt,
            headheight=14pt,headsep=12pt,footskip=22pt]{geometry}
\usepackage{xcolor}
\usepackage{microtype}
\usepackage{ragged2e}
\usepackage{booktabs}
\usepackage{longtable}
\usepackage{tabularx}
\usepackage{array}
\usepackage{enumitem}
\usepackage{needspace}
\usepackage{fancyhdr}
\usepackage{hyperref}

\definecolor{molariumblue}{HTML}{14607E}
\definecolor{runninggray}{gray}{0.50}

\hypersetup{
  pdftitle={Molarium Appendix B - Scientific workflows for humans and agents},
  pdfauthor={Rafal P. Wiewiora, W. Patrick Walters, and Woody Sherman},
  pdfsubject={Task-centered operating appendix for Molarium},
  colorlinks=false,
  pdfborder={0 0 0}
}

\pagestyle{fancy}
\fancyhf{}
\renewcommand{\headrulewidth}{0.4pt}
\fancyhead[L]{\color{runninggray}\fontsize{8.97}{10.5}\selectfont Molarium - Appendix B}
\fancyhead[R]{\color{runninggray}\fontsize{8.97}{10.5}\selectfont Wiewiora, Walters, and Sherman}
\fancyfoot[C]{\fontsize{9.96}{11.95}\selectfont\thepage}

\setcounter{page}{19}
\setlength{\parindent}{0pt}
\setlength{\parskip}{5.5pt}
\setlength{\emergencystretch}{2em}
\renewcommand{\arraystretch}{1.08}
\AtBeginDocument{\fontsize{9.96}{11.95}\selectfont\justifying}

\setlist[itemize]{leftmargin=14pt,labelsep=5pt,itemsep=1pt,topsep=2pt,parsep=0pt}
\setlist[enumerate]{leftmargin=18pt,labelsep=5pt,itemsep=1.5pt,topsep=2pt,parsep=0pt}

\newcommand{\code}[1]{\texttt{#1}}
\newcommand{\lead}[1]{\textbf{#1}\hspace{0.35em}}
\newcommand{\subhead}[1]{%
  \par\needspace{2\baselineskip}\vspace{2.5pt}%
  {\fontsize{10.91}{13.1}\selectfont\bfseries #1}\par\vspace{-1.5pt}}

\newcommand{\workflowsection}[4]{%
  \par\needspace{5\baselineskip}\vspace{7.5pt}%
  {\fontsize{10.91}{13.1}\selectfont\bfseries #1\hspace{1.1em}#2}\par\vspace{2pt}%
  \lead{Status.}#3.\par\vspace{-3.5pt}%
  \lead{The question.}#4\par}

\newcolumntype{K}{>{\bfseries\RaggedRight\arraybackslash}p{1.24in}}
\newcolumntype{V}{>{\RaggedRight\arraybackslash}p{5.53in}}
\newenvironment{skillcontract}[1]{%
  \par\needspace{5\baselineskip}\vspace{2pt}%
  {\fontsize{10.91}{13.1}\selectfont\bfseries Agent skill contract: \code{#1}}\par\vspace{1pt}%
  \fontsize{8.4}{10}\selectfont
  \begin{longtable}{@{}K V@{}}
  \toprule
  \textbf{Contract element} & \textbf{Requirement}\\
  \midrule
  \endfirsthead
  \toprule
  \textbf{Contract element} & \textbf{Requirement}\\
  \midrule
  \endhead
  \midrule
  \multicolumn{2}{r}{\footnotesize Continued on next page}\\
  \endfoot
  \bottomrule
  \endlastfoot
}{%
  \end{longtable}
  \normalsize\fontsize{9.96}{11.95}\selectfont}
\newcommand{\contractrow}[2]{#1 & #2\\}

\newenvironment{smallcomparison}[1]{%
  \par\vspace{2pt}\fontsize{8.4}{10}\selectfont
  \begin{longtable}{@{}#1@{}}\toprule
}{%
  \bottomrule\end{longtable}
  \normalsize\fontsize{9.96}{11.95}\selectfont}

\begin{document}

\begin{center}
{\color{molariumblue}\fontsize{14.35}{17}\selectfont\bfseries APPENDIX B\par}
\vspace{2pt}
{\color{molariumblue}\fontsize{17.22}{20.5}\selectfont\bfseries
Working with Molarium: scientific workflows for humans and agents\par}
\end{center}
\vspace{11pt}

Molarium is designed so that a scientist working through the graphical interface and an agent operating through the Chemist Actions API act on the same molecular state and encounter the same scientific boundaries. This appendix therefore describes Molarium by \textbf{scientific task}, not by software module.

Each workflow begins with the question a scientist is trying to answer. The human procedure explains how to perform the task and what to inspect before accepting the result. The accompanying \textbf{agent skill contract} expresses the same procedure as preconditions, permitted actions, evidence, success criteria, and stopping conditions. The two views are intentionally paired: a procedure should be understandable by a scientist and sufficiently explicit that an agent cannot silently change its meaning.

\lead{Implementation basis.}This appendix describes the inspected browser implementation on 2 September 2026, \code{feature/sos1-interface-story} through commit \code{ae052c4}. It distinguishes current browser behavior from executable but experimental behavior. Repository-only and proposed functionality is not presented as an operating instruction.

\workflowsection{B.1}{One molecular state, two interfaces}{Current browser surface}{What changes the molecule, and what merely changes how I am looking at it?}

Molarium separates the authoritative molecular state from the interface used to manipulate it. A scientist can work in \textbf{View}, \textbf{Build}, and \textbf{Run}; an agent can invoke the corresponding guarded operations through the versioned Chemist Actions API. Both routes act on the same graph, coordinates, selections, calculation records, and pending chemistry transaction.

Rotating the camera, changing a representation, focusing a component, showing contacts, or selecting a saved frame changes the view. Loading a structure, applying a protonation state, finishing a chemistry transaction, accepting a docking candidate, or applying optimized coordinates changes the authoritative molecular state. Candidate structures and poses remain candidates until explicitly applied.

\lead{The common operating rule.}Inspect, act, verify. Generate alternatives freely, but change the authoritative molecular state deliberately. A valid operation is not automatically a scientifically appropriate one; acceptance still depends on the question being asked.

\subhead{Starting from a structure}

Start with \code{?blank} when a reproducible workflow should begin from an empty canvas. Otherwise the application loads its bundled launch molecule. A structure can be pasted as XYZ or PDB text, uploaded as XYZ or PDB, entered as a SMILES string or PDB identifier in connected mode, or selected from the prepared examples. The current upload router should not be relied on for ordinary MOL-file ingestion, even though a V2000 parser exists elsewhere in the application.

After loading, inspect atom count, formula, formal charge, connected components, graph connectivity, and the 3D structure. For a selected small-molecule component, compare the synchronized 2D depiction with the 3D graph. For PDB input, remember that only the first model is loaded and alternate locations are resolved by the implemented occupancy policy. For XYZ input, bonds are inferred from proximity and should be checked rather than assumed.

\subhead{State and persistence}

Loading a new molecule resets incompatible calculation, docking, selection, and preparation state. The working session is in memory: Clear, page reload, or script import can discard unsaved work. Share copies a URL rather than serializing the current molecule, and Save produces a PNG rather than a reusable molecular record. Export the scientific record needed by the next step before replacing or clearing the scene.

\begin{skillcontract}{operate\_molecular\_state}
\contractrow{Goal}{Load and inspect a starting structure without confusing a visual change, candidate, or imported approximation with an accepted molecular state.}
\contractrow{Preconditions}{A supported input or prepared example is available; connected-mode requirements are known.}
\contractrow{Inputs}{Input type and contents; intended component; optional network policy; requested inspection fields.}
\contractrow{Procedure}{Start blank when reproducibility requires it; load through a supported route; inspect graph, components, charge, and coordinates; report ambiguities before any mutation.}
\contractrow{State mutation}{A successful load replaces the current molecule and resets dependent state. View and selection operations do not alter graph chemistry.}
\contractrow{Evidence}{Source identity or input hash where available; atom/component summary; explicit warnings about inferred bonds, alternate locations, or network retrieval.}
\contractrow{Success}{The loaded graph and coordinates correspond to the intended starting object and all material ambiguities have been surfaced.}
\contractrow{On failure}{Leave the prior scene in place where the interface does so; report the parse, fetch, or format error; do not reinterpret an unsupported file silently.}
\contractrow{Boundary}{No autosave or complete-session archive; an image or share link is not a reusable molecular state.}
\end{skillcontract}

\workflowsection{B.2}{Prepare a protein for local molecular design}{Experimental current surface}{I have an experimental protein structure. What must happen before I can safely edit a ligand or run a local calculation?}

A deposited PDB structure is not automatically a simulation-ready system. Atoms may be missing, alternate locations may have been chosen, protonation can be ambiguous, crystallographic waters may or may not matter, and ligands or unusual residues may exceed the supported parameterization. Molarium's preparation workflow is deliberately conservative: it makes a bounded set of assumptions explicit, performs repairs it knows how to perform, and stops when the structure exceeds those capabilities.

\subhead{Human workflow}

\begin{enumerate}
\item Load the PDB structure and inspect the protein, ligand, waters, other components, and the region around the intended binding site before preparing anything.
\item Open \textbf{Protein Preparation}. Choose pH, histidine policy, supported missing-heavy-atom repair, ligand treatment, water treatment, and gap policy.
\item Generate the preparation preview. It inventories residues, missing atoms, sequence gaps, ligands, waters, clashes, and invalid coordinates. Within scope, it can retrieve CCD information, reconstruct supported missing side-chain atoms, add standard hydrogens and terminal oxygen, choose histidine states, and scan polar-hydrogen orientations.
\item Inspect every proposed change, warning, and blocker. Download the preparation report before applying if the decisions must survive the browser session.
\item Apply only when the declared preparation is appropriate for the intended bounded calculation. Apply replaces the graph and coordinates with the prepared state and attaches the preparation audit and numerical parameterization.
\end{enumerate}

\subhead{What should I inspect?}

\begin{itemize}
\item Ligand identity, connectivity, completeness, and whether it should be retained or excluded.
\item Histidine and other local protonation choices near the ligand or catalytic region.
\item Crystallographic waters that bridge interactions or define the local geometry.
\item Missing atoms, chain breaks, unresolved regions, reconstructed geometry, and severe clashes near the region of interest.
\item Whether the resulting force-field and charge assumptions match the scientific question.
\end{itemize}

\lead{Acceptance is narrower than success.}A clean preview means Molarium completed its declared preparation procedure. It does not certify a production molecular-dynamics system. The implementation explicitly records \code{readyForProductionSimulation: false}.

Preparation blocks unsupported residues, unsupported missing atoms, incomplete retained ligands, disallowed gaps, invalid coordinates, and severe clashes. It does not currently construct missing loops or membranes, model general metal coordination, parameterize covalent ligands, or build periodic explicit-solvent boxes. Arbitrary proteins use an experimental Sage/Gasteiger route; exact prepared fixtures do not demonstrate general browser-native protein typing.

\begin{skillcontract}{prepare\_protein}
\contractrow{Goal}{Convert the loaded PDB into an explicitly reviewed molecular state suitable for a supported local Molarium calculation.}
\contractrow{Preconditions}{A PDB-derived structure is loaded; the intended downstream use is known; ligand and water roles can be reviewed.}
\contractrow{Inputs}{pH; histidine policy; missing-heavy-atom repair policy; ligand policy; water policy; gap policy.}
\contractrow{Procedure}{Inspect starting structure; generate preview; review issues and proposed changes; export the report when persistence matters; stop on unresolved blockers; apply only an accepted preparation.}
\contractrow{State mutation}{Prepared graph and coordinates replace the prior state; preparation audit and parameterization are attached; one undo snapshot is created.}
\contractrow{Evidence}{Options, detected issues, repairs, warnings and blockers, preparation report, and resulting parameterization identity.}
\contractrow{Success}{The declared procedure completes without unresolved blockers and its scope matches the intended local calculation.}
\contractrow{Failure behavior}{Report the blocker and stop. Do not manufacture missing atoms, parameters, loop geometry, or unsupported chemistry.}
\contractrow{Escalate when}{The problem requires loops, membranes, metal coordination, covalent chemistry, periodic solvent, or another unsupported preparation capability.}
\contractrow{Does not establish}{Production simulation readiness or physical accuracy of the chosen downstream model.}
\end{skillcontract}

\workflowsection{B.3}{Choose the protonation state you are designing}{Current ligand workflow; protein choices within experimental preparation}{What charged and protonated species does the molecular graph actually represent?}

Protonation is part of chemical identity, not a cosmetic display choice. It changes formal charge, hydrogen count, hydrogen-bonding patterns, parameterization, and often the meaning of every calculation that follows. Molarium exposes two related but distinct workflows: ligand-state enumeration through RDKit, and protein protonation decisions inside the bounded PDB-preparation workflow.

\subhead{Ligand workflow}

\begin{enumerate}
\item Load the ligand from a SMILES string or select the ligand component whose identity you intend to change.
\item Choose a pH from 0 to 14 and enumerate candidate states. RDKit applies a versioned Dimorphite-style site table, edits formal charges and hydrogens, sanitizes each result, and returns an ordered list with estimated weights.
\item Inspect each candidate as chemistry: total charge, changed atoms, tautomeric or aromatic consequences, and compatibility with known assay conditions or the bound environment.
\item Apply one state deliberately. Molarium re-embeds the chosen species, replaces the graph and coordinates, and invalidates state that depended on the former topology.
\end{enumerate}

\subhead{Protein workflow}

For a PDB structure, set the pH and histidine policy during preparation. The implemented policies include automatic selection or explicit HID, HIE, and HIP choices, followed by inspection of polar-hydrogen orientations. Treat residues near a ligand, metal, catalytic center, or salt-bridge network as scientific decisions rather than defaults.

\lead{What the reported ligand weights mean.}They are independent Henderson-Hasselbalch estimates from the implemented site model. They are not microscopic pKa predictions, coupled-site populations, or evidence for the state selected by a protein binding site.

\begin{skillcontract}{select\_protonation\_state}
\contractrow{Goal}{Make the intended protonation and formal-charge state explicit before editing or calculation.}
\contractrow{Preconditions}{A valid ligand or PDB structure is loaded; the relevant pH or experimental condition is known or stated as an assumption.}
\contractrow{Inputs}{pH; candidate-state selection; for proteins, histidine policy and local residues requiring review.}
\contractrow{Procedure}{Enumerate or preview; compare candidates; inspect charge and local interactions; apply exactly one accepted state; re-inspect the resulting graph.}
\contractrow{State mutation}{Applying a ligand state replaces graph and coordinates. Protein protonation changes are committed through preparation. Prior parameterization becomes obsolete.}
\contractrow{Evidence}{pH, enumeration or preparation output, chosen state, total charge, changed sites, and stated rationale when more than one state is plausible.}
\contractrow{Success}{The authoritative graph explicitly represents the intended state and downstream methods will see the corresponding atoms and charges.}
\contractrow{On failure}{Report invalid pH, sanitization, embedding, or worker errors; do not partially apply or choose a different state merely because it runs.}
\contractrow{Boundary}{Enumeration weights do not resolve coupled microscopic states or protein-induced pKa shifts.}
\end{skillcontract}

\workflowsection{B.4}{Edit atoms, bonds, charges, and stereochemistry in 3D}{Current browser surface}{How do I change the chemistry without corrupting the molecular graph or losing the ability to undo the operation?}

The synchronized 2D and 3D views are two ways of selecting the same graph. Chemistry-changing actions are grouped into a transaction so that a related set of edits can be validated and committed as one scientific step. Protein covalent atoms are protected from ordinary ligand editing.

\subhead{Human workflow}

\begin{enumerate}
\item Enter \textbf{Build} and select the relevant atom or bond in either view. Confirm the stable atom identity and local environment before changing it.
\item Apply the intended element replacement, bond order or aromaticity change, formal charge, explicit hydrogen, or stereochemical operation. Related operations remain in one pending chemistry transaction.
\item For coordinate-only geometry changes, select a connected path of two, three, or four atoms to set a distance, angle, or dihedral. Molarium moves the connected side defined by the graph and blocks ambiguous ring cuts.
\item Choose \textbf{Finish chemistry}. RDKit sanitizes the graph, hydrogens are reconciled, eligible local coordinates are polished, reference contacts are remapped where applicable, obsolete parameterization is invalidated, and a single undo snapshot is created.
\item If validation fails, correct the still-pending transaction or choose \textbf{Discard} to restore the pre-edit state. Use Undo/Redo only after a transaction has been committed.
\end{enumerate}

\subhead{What should I inspect?}

\begin{itemize}
\item Valence, formal charge, aromaticity, hydrogen count, stereochemistry, and the intended bond topology.
\item Whether the 2D depiction and 3D graph agree after sanitization.
\item New close contacts, distorted local geometry, and movement outside the edited region.
\item Whether topology-dependent records---parameterization, docking mapping, or calculation results---were invalidated or remapped as expected.
\end{itemize}

A synchronous mutation error restores the transaction snapshot. If the chemistry is valid but optional coordinate polishing fails, the valid graph can remain committed and the cleanup error is recorded. Fused-aromatic edits that cannot preserve the aromatic graph are rejected rather than approximated.

\begin{skillcontract}{edit\_molecular\_graph}
\contractrow{Goal}{Apply an explicit, reviewable graph or local-geometry change while preserving atomic lineage and undoability.}
\contractrow{Preconditions}{The intended editable component is identified; no unrelated transaction is pending; the target atoms or bond are unambiguous.}
\contractrow{Inputs}{Selected stable atom identifiers; edit type and value; optional connected geometry path; intended transaction boundary.}
\contractrow{Procedure}{Inspect target; open or join a transaction; apply supported edits; finish once; inspect sanitized graph and local coordinates.}
\contractrow{State mutation}{Graph edits remain provisional until Finish. Finish commits one new state and invalidates obsolete parameterization; Discard restores the starting state.}
\contractrow{Evidence}{Pre-edit identity, ordered public actions, changed atoms and bonds, sanitization outcome, cleanup outcome, and final inspection.}
\contractrow{Success}{The final graph encodes the intended chemistry, sanitizes successfully, and passes local visual and structural checks.}
\contractrow{On failure}{Keep the transaction pending for correction or discard it. Never bypass sanitization or inject coordinates directly.}
\contractrow{Boundary}{Local graph validity does not establish synthetic feasibility, tautomer preference, binding affinity, or global conformational plausibility.}
\end{skillcontract}

\workflowsection{B.5}{Grow an analogue by adding and positioning fragments}{Current builder; pose-aware extensions experimental}{How do I add the substituent I have in mind while preserving the useful context of the parent structure?}

Fragment growth combines a graph edit with an initial 3D placement. Those are separate scientific problems: the new molecule must first be chemically correct, and its coordinates must then be a plausible starting point for relaxation or pose maintenance. Molarium's Builder stages fragment templates and commits the resulting chemistry through the same transaction system used for simpler edits.

\subhead{Human workflow}

\begin{enumerate}
\item If the parent is bound and its pose matters, capture the reference pose \textbf{before} beginning the topology change. Retain any explicit core or interaction choices needed for later pose maintenance.
\item Select the ligand attachment atom and stage the desired fragment or template. Confirm which atoms are new, which parent atom remains the anchor, and which bond is to be formed.
\item Create the attachment and perform any associated leaving-group, bond-order, formal-charge, hydrogen, or stereochemical edits within the same transaction.
\item Use distance, angle, and dihedral controls to give the new group a sensible initial orientation. These controls establish a starting geometry; they do not score the binding site.
\item Finish the chemistry, then inspect the attachment, local valence, stereochemistry, ring geometry, clashes, and preserved parent coordinates. Relax the structure or invoke reference-guided pose maintenance as a separate step.
\end{enumerate}

\subhead{How should I judge the result?}

A successful fragment operation means that the intended graph was built and sanitized. It does not mean that the chosen rotamer, ring conformation, or pocket orientation is preferred. For a solvent-exposed edit, local minimization may be enough to remove construction strain. For an edit that changes pocket contacts or releases a ring, use a workflow that samples alternatives and retains the relationship to the experimental reference.

\lead{Enumeration boundary.}Development modules can represent bounded edit plans, but general combinatorial enumeration, library management, automated campaign selection, and automatic MCS transfer of an independently imported analogue are not current workbench workflows.

\begin{skillcontract}{grow\_fragment}
\contractrow{Goal}{Construct one intended analogue and a reviewable starting geometry while retaining parent atom lineage and, when requested, reference-pose context.}
\contractrow{Preconditions}{Editable ligand and attachment atom identified; fragment chemistry specified; reference captured first if pose continuity matters.}
\contractrow{Inputs}{Parent and attachment atom identifiers; fragment/template; new bond; associated deletions, charge or stereochemical edits; optional geometry targets.}
\contractrow{Procedure}{Capture reference if needed; stage fragment; perform one coherent transaction; set only the geometry needed for a starting pose; finish and inspect; relax or sample separately.}
\contractrow{State mutation}{Finish commits the new graph, reconciles hydrogens, and invalidates prior parameterization. A later pose or optimization changes coordinates separately.}
\contractrow{Evidence}{Parent and product inspection, atom mapping, edit actions, transaction outcome, initial placement rationale, and any preserved reference/contact record.}
\contractrow{Success}{The analogue graph is correct and sanitized; the initial geometry has no obvious construction pathology and is suitable for the chosen next workflow.}
\contractrow{On failure}{Correct or discard the pending transaction. Do not treat failed cleanup as permission to accept an invalid graph or arbitrary pose.}
\contractrow{Boundary}{Fragment growth is not combinatorial design, synthesis planning, global docking, or proof that the new group is energetically favored.}
\end{skillcontract}

\workflowsection{B.6}{Relax a structure, run bounded dynamics, and inspect motion}{Current and experimental methods with different scopes}{I have made a structure. Does its geometry survive relaxation, and what moves during a short calculation?}

Choose the method from the scientific question, not from a generic expectation that every engine answers the same thing. Molarium exposes several numerical pathways whose energies are not interchangeable. Record the selected model, solvent, constraints, and protocol before interpreting the result.

\begingroup
\fontsize{7.8}{9.2}\selectfont
\begin{longtable}{@{}p{1.31in}p{0.78in}p{1.77in}p{1.85in}@{}}
\toprule
\textbf{Question} & \textbf{Method} & \textbf{What it provides} & \textbf{Principal boundary}\\
\midrule
\endfirsthead
\toprule
\textbf{Question} & \textbf{Method} & \textbf{What it provides} & \textbf{Principal boundary}\\
\midrule
\endhead
Small-molecule cleanup or short force-field motion & RDKit & Energy, geometry optimization, and short dynamics using the RDKit force-field path. & The reported 0.1 fs dynamics path is a workbench calculation, not a production trajectory.\\
Sage energy, relaxation, or bounded single-system MD & Direct WebGPU & Browser-GPU energy, forces, geometry, and saved frames from the numeric Sage system. & Experimental; f32; direct dynamics capped at 5,000 steps; no PBC or PME.\\
Many replicas of one topology & WebGPU ensemble & Homogeneous replica dynamics with energies, frames, timing, and constraint diagnostics. & Experimental; at most 512 atoms per replica, 1,024 replicas, and 100,000 steps; nonperiodic all-pairs model.\\
Electronic energy or geometry for an eligible neutral molecule & ANI-2x & Eight-model ANI-2x ensemble energy, spread, analytical forces, and optional optimized coordinates. & Experimental; connected neutral singlets with explicit H; H/C/N/O/S/F/Cl only; at most 96 atoms.\\
Reference calculation inside another workflow & OpenMM 8.2 Reference & Double-precision internal parameterization, scoring, relaxation, and validation path. & Internal subsystem, not a selectable general Run method.\\
\bottomrule
\end{longtable}
\endgroup

\subhead{Human workflow}

\begin{enumerate}
\item Finish any pending chemistry. Confirm that the molecule is eligible for the chosen engine and, where required, can be completely parameterized.
\item In \textbf{Run}, choose the job and compatible method. Set the step count, vacuum or OBC2/ACE implicit solvent, X-H constraints where supported, saved frames, and ensemble options.
\item Run once. Inspect the engine or model identity, force field, charge model, solvent, constraints, initial and final energy, convergence or diagnostic status, elapsed time, and all warnings.
\item For dynamics, step through saved frames or replicas and inspect the region relevant to the design question. Display alignment does not overwrite raw calculation coordinates.
\item Apply or export coordinates only when the result passes the declared acceptance check. Current-frame XYZ is available; saved frames are a bounded sample, and the workbench has no generic full-trajectory-file export.
\end{enumerate}

\subhead{Interpreting the model}

Automatic Sage parameterization uses deterministic Gasteiger charges rather than AM1-BCC or NAGL. Unsupported elements, invalid chemistry, or missing required parameters fail instead of silently dropping terms. The exposed workflows are nonperiodic and do not provide PME, a barostat, explicit-solvent setup, generic production protein simulation, or binding free energy.

\begin{skillcontract}{relax\_or\_simulate}
\contractrow{Goal}{Use a declared numerical model to test local geometry or bounded motion and retain enough metadata to interpret the result.}
\contractrow{Preconditions}{No chemistry transaction is pending; method eligibility and complete parameterization are established; an acceptance observable is specified.}
\contractrow{Inputs}{Job; method; steps; solvent; constraints; saved frames; fixed or movable atoms where supported; ensemble settings.}
\contractrow{Procedure}{Validate eligibility; record model and protocol; run one compatible job; inspect energies, diagnostics, and frames; apply coordinates only after acceptance.}
\contractrow{State mutation}{Energy jobs do not move the molecule. Successful geometry or dynamics jobs can replace current coordinates; raw frames remain calculation records.}
\contractrow{Evidence}{Engine/model version and hashes where recorded, force field, charge model, solvent, constraints, timestep, energies, convergence/diagnostics, timing, frames, and warnings.}
\contractrow{Success}{The worker returns a finite, internally valid result and the predefined structural check is satisfied within the method's declared scope.}
\contractrow{On failure}{Keep the prior graph and coordinates; record the returned error; correct eligibility, options, or assets before retrying.}
\contractrow{Boundary}{A minimized pose or short stable trace is not an affinity estimate, equilibrium ensemble, free-energy calculation, or production trajectory.}
\end{skillcontract}

\workflowsection{B.7}{Explore conformations with Conformer Arena}{Experimental current surface}{What other conformations can this molecule adopt, and how sensitive is the answer to the generation and refinement method?}

Conformer search is an alternative-generation problem. Molarium separates seed generation, physical refinement, clustering, comparison, and selection so that a low-energy-looking structure is not confused with proof that the ensemble is complete or Boltzmann weighted.

\subhead{Human workflow}

\begin{enumerate}
\item Start from a chemically finished small molecule with explicit hydrogens and a sensible protonation state. Recenter the structure before tight STORMM positional comparisons, because fixed-point spatial precision degrades far from the working origin.
\item Choose \textbf{WebGPU ensemble} and \textbf{Conformer search}. Set the requested conformer count, effort, clustering cutoff, solvent/constraint options, and whether to open the comparison arena.
\item Molarium generates deterministic RDKit conformer seeds, up to the implemented cap of 256. The STORMM protocol then performs staged relaxation, heating at 600 K, cooling through 450 K and 300 K, and final relaxation across homogeneous replicas.
\item Inspect the returned ensemble before selecting a representative: energy range, distinct RMSD clusters, cluster populations, recurring torsions, ring states, constraint diagnostics, and obvious duplicates or failures.
\item Use Conformer Arena to compare supported pathways on the same molecule and settings. Ask whether the methods discover the same basins, not merely whether they rank one shared conformer similarly.
\item Apply a selected conformer only after inspection. Export the ensemble as SDF when it fits V2000 record limits; retain method and clustering metadata separately.
\end{enumerate}

\subhead{What does the result establish?}

The workflow shows which basins the declared finite search found and how the supported methods behaved on that search. The anneal/minimize protocol is a search heuristic, so replica counts and cluster frequencies are not thermodynamic populations. Failure to find a conformation is not evidence that it is inaccessible, and energy differences from different force fields are not directly comparable.

\begin{skillcontract}{explore\_conformers}
\contractrow{Goal}{Generate, refine, cluster, and compare a finite set of conformational alternatives for one molecular topology.}
\contractrow{Preconditions}{Chemistry and protonation are finalized; molecule is eligible; intended use and comparison metric are stated.}
\contractrow{Inputs}{Conformer count; effort; cluster cutoff; seed; solvent and constraints; arena choice; acceptance or diversity criteria.}
\contractrow{Procedure}{Generate RDKit seeds; run declared refinement protocol; cluster and inspect; compare basins in Arena when requested; select or export without rewriting raw records.}
\contractrow{State mutation}{Search produces a candidate ensemble. Only explicit selection changes the displayed/current coordinates.}
\contractrow{Evidence}{Seed and method identity, requested and returned counts, energies, cluster assignments, diagnostics, timings, and exported ensemble where applicable.}
\contractrow{Success}{A finite set of valid, distinct candidates is produced and its diversity is sufficient for the stated downstream purpose.}
\contractrow{On failure}{Report parameterization, eligibility, replica, overflow, or pair-work limits; reduce scope or escalate rather than hiding missing replicas.}
\contractrow{Boundary}{Not exhaustive conformational sampling, a Boltzmann ensemble, solution free energies, or proof of bioactive-conformation probability.}
\end{skillcontract}

\workflowsection{B.8}{Predict a short single-chain protein structure}{Experimental current surface; connected mode only}{Can Molarium create a protein structure from sequence, and is this the co-folding workflow?}

\lead{This is not currently co-folding.}The implemented workflow predicts one protein entity from one sequence. It does not accept a ligand for joint protein-ligand prediction, does not predict a multimer, and does not maintain a starting ligand pose. Calling it co-folding would overstate the present capability.

The Fold Protein panel combines a configured remote MSA service with local OpenFold ONNX inference. The sequence is transmitted to the selected MSA endpoint; model inference then runs in the browser, preferring WebGPU and falling back to WASM. Local Lab disables the external MSA path.

\subhead{Human workflow}

\begin{enumerate}
\item Enter one amino-acid sequence using the standard residue alphabet or X. The current maximum length is 128 residues, using 64- or 128-residue feature buckets.
\item Inspect and accept the configured MSA endpoint and network boundary. First use may require approximately 540 MB of model assets.
\item Choose recycle settings; the current default is three. Run and monitor MSA and inference progress. Cancel aborts the active MSA/prediction controller.
\item On success, Molarium loads the prediction as the current molecule. Inspect chain completeness, residue mapping, gross geometry, and suitability for the downstream question before replacing other work.
\item Export the predicted structure as PDB if it must persist. Preparation and local relaxation are separate decisions.
\end{enumerate}

The current path does not use templates or perform Amber relaxation. It creates no server-side job ledger. Fetch, archive, feature, model, or inference failure produces no partial prediction. A returned structure is a model output, not experimental evidence or a prepared simulation system.

\begin{skillcontract}{predict\_protein\_structure}
\contractrow{Goal}{Predict coordinates for one short protein sequence using the declared MSA service and local OpenFold model.}
\contractrow{Preconditions}{Connected mode; configured MSA endpoint; sequence length at most 128; network and model-asset use accepted.}
\contractrow{Inputs}{One sequence; MSA endpoint; recycle settings.}
\contractrow{Procedure}{Validate sequence; disclose external transmission; obtain and validate MSA; run the matching ONNX bucket; inspect result; export before replacing it.}
\contractrow{State mutation}{A successful prediction replaces the current molecule. Failure does not load a partial structure.}
\contractrow{Evidence}{Input sequence and endpoint, model/bucket and recycle settings, progress/failure record, prediction metadata, and exported PDB where required.}
\contractrow{Success}{A complete finite prediction is returned for the declared input and passes basic structural inspection for the intended next step.}
\contractrow{On failure}{Report MSA, archive, rate-limit, maintenance, timeout, feature, asset, or inference failure; do not fabricate a partial structure.}
\contractrow{Boundary}{No ligand co-folding, multimer prediction, templates, Amber relaxation, experimental validation, or production preparation.}
\end{skillcontract}

\workflowsection{B.9}{Maintain an experimental pose while changing ligand chemistry}{Experimental current surface}{I trust the experimental pose of my starting ligand. How can I search for a plausible pose of an edited analogue without discarding what I already know?}

Reference-guided docking is a local analogue pose-maintenance workflow. It preserves stable atom identities and selected interactions from a prepared reference complex, searches the degrees of freedom introduced by an edit, rejects infeasible candidates, and ranks distinct feasible poses in a rigid local receptor site. It is intentionally narrower than global docking.

\subhead{Human workflow}

\begin{enumerate}
\item Prepare and inspect the reference protein-ligand complex. In \textbf{Build}, capture the reference before editing. The capture stores ligand lineage, receptor atoms within 8~\AA, reference contacts, selected hydrogen bonds, settings, and provenance.
\item Use the default \textbf{Pose propagation} mode to map surviving reference atoms and seed the edited region. Use \textbf{Expert selected core} only when you can justify at least three connected, non-collinear heavy atoms as a fixed reference.
\item Edit and finish the analogue. Review remapped or unavailable contacts, transformed ring regions, atom mapping, and any constraints before starting the search.
\item Run the constrained search. Molarium verifies receptor and contact integrity, releases eligible transformed rings, generates deterministic feature and torsion seeds, freshly parameterizes the edited ligand, relaxes feasible candidates, and scores the distinct results.
\item Inspect the ranked candidates, not only rank 1. Compare mapped-core displacement, preserved or replaced interactions, ligand strain, clashes, ring and side-chain choices, and the reasons candidates were rejected.
\item Apply one candidate explicitly. Candidate generation does not change the authoritative molecule. Application rechecks topology and receptor integrity before replacing coordinates. Export the labbook as JSON and human-readable Markdown/notes.
\end{enumerate}

\subhead{What is being scored?}

The ranker uses a rigid receptor site within 8~\AA, Lennard-Jones and Coulomb ligand-receptor cross terms with relative dielectric 4, and relative vacuum ligand strain. It omits general receptor relaxation, water displacement, entropy, binding free energy, and broad macrocycle or fused-ring concerted search. The separate induced-fit Builder operation is not part of this ranker and is not an affinity calculation.

\lead{Negative results are results.}An invalid core, altered receptor, unmappable required atom, unavailable contact feature, parameterization failure, or absence of feasible candidates produces an explicit failed or negative record. Molarium does not silently apply the least-bad pose.

\begin{skillcontract}{maintain\_reference\_pose}
\contractrow{Goal}{Generate and review plausible local poses for one edited analogue while retaining explicit relationships to a trusted reference complex.}
\contractrow{Preconditions}{Prepared complex; trusted reference pose; reference captured before edit; chemistry finished; receptor unchanged; required contacts available.}
\contractrow{Inputs}{Pose-propagation or expert-core mode; selected core or contacts; cleanup/contact policy; candidate count; analogue graph.}
\contractrow{Procedure}{Capture; edit; finish; verify mapping and receptor integrity; generate and score feasible candidates; inspect alternatives; apply one explicitly or retain a negative result.}
\contractrow{State mutation}{Capture and search create reference and candidate records. Only Apply replaces the current coordinates.}
\contractrow{Evidence}{Protocol/settings snapshot, hashes, atom mapping, contact decisions, constraints, rejected candidates, scores, selected pose, and chained labbook.}
\contractrow{Success}{At least one feasible distinct candidate survives the declared constraints and a human or agent can justify any applied selection from the recorded evidence.}
\contractrow{On failure}{Apply nothing; preserve the negative labbook; repair the reference, mapping, constraints, or chemistry before rerunning.}
\contractrow{Boundary}{Local analogue pose maintenance only---not global docking, automatic MCS pose transfer, receptor-ensemble docking, or affinity/free-energy prediction.}
\end{skillcontract}

\workflowsection{B.10}{Record, replay, and hand off a design process}{Current action/replay surface; campaign ledger separate and experimental}{I have done something scientifically useful. What must I export so that another scientist or an agent can understand and reproduce it?}

Molarium records several different kinds of evidence. They answer different questions and should not be treated as interchangeable. The live action audit says what was requested; a calculation or workflow record says what method ran and what it returned; an action script says how to replay completed public operations; a campaign ledger represents a durable sequence of content-addressed decisions. The current workbench does not automatically combine them into one project archive.

\begingroup
\fontsize{8.0}{9.5}\selectfont
\begin{longtable}{@{}p{1.23in}p{2.76in}p{2.61in}@{}}
\toprule
\textbf{Record} & \textbf{Question answered} & \textbf{Persistence and limit}\\
\midrule
\endfirsthead
\toprule
\textbf{Record} & \textbf{Question answered} & \textbf{Persistence and limit}\\
\midrule
\endhead
Chemist Actions audit & What recognized public action was requested, with what arguments and outcome? & Session memory; serialized calls; bounded history, normally 500 entries; cleared with the scene unless exported through a script path.\\
Preparation or calculation record & What system, model, settings, energies, frames, warnings, or audit produced this result? & Attached to the in-memory state and available through workflow-specific reports or exports.\\
Docking labbook & How were mapping, feasibility, scoring, rejection, and selection handled? & Exportable JSON and Markdown/notes with chained hashes.\\
Designer action script & Which completed public actions should be replayed? & Exportable JSON; replay invokes the same guarded routes; not a complete UI or binary session snapshot.\\
Registered design route & What prospective, hash-pinned edit/evaluation sequence is permitted? & Validated route schema; distinct from retrospective history.\\
Campaign ledger & What content-addressed structures, decisions, commits, and branches form the design history? & Experimental/repository surface with append-only and finalized records; not one-click promotion of every live session.\\
\bottomrule
\end{longtable}
\endgroup

\subhead{Replay workflow}

\begin{enumerate}
\item Export any unsaved molecular or workflow records before importing a script; import intentionally clears the existing scene.
\item Create a Designer action script from completed recognized public actions. Failed actions remain audit evidence but are not silently converted into successful replay steps.
\item Import the validated script, restart it, and Play or step through it. Replay uses the same Chemist Actions routes, validations, and workers as manual operation and produces a separate replay audit.
\item Compare the resulting graph, coordinates, decisions, and errors with the intended endpoint. Export the replay audit and scientific records needed by the recipient.
\item For a registered story or route, verify schema, referenced content, and SHA-256 before execution. Treat a hash chain as integrity evidence, not an authenticated author identity or proof of scientific truth.
\end{enumerate}

\subhead{Choose the export that matches the handoff}

\begingroup
\fontsize{8.0}{9.5}\selectfont
\begin{longtable}{@{}p{1.20in}p{2.48in}p{2.92in}@{}}
\toprule
\textbf{Export} & \textbf{Carries} & \textbf{Does not carry}\\
\midrule
\endfirsthead
\toprule
\textbf{Export} & \textbf{Carries} & \textbf{Does not carry}\\
\midrule
\endhead
XYZ current frame & Displayed coordinates & Full topology metadata, history, or provenance\\
SDF conformer ensemble & Exportable conformer graph and coordinates & Unrelated session state or complete search history\\
PDB protein & Imported or predicted protein coordinates and supported annotations & Complete application or preparation provenance\\
Preparation JSON & Policies, issues, audit, and readiness status & A production-readiness guarantee\\
Docking JSON/Markdown & Settings, mapping, candidates, failures, hashes, and selection & Author authentication or binding free energy\\
Action script JSON & Public actions and bindings & Arbitrary code, direct coordinates, or complete view state\\
PNG or share URL & A visual snapshot or current/registered URL & A reusable molecule or serialized live session\\
\bottomrule
\end{longtable}
\endgroup

\begin{skillcontract}{export\_and\_handoff}
\contractrow{Goal}{Preserve the smallest complete set of molecular, procedural, numerical, and decision records needed by the next scientist or agent.}
\contractrow{Preconditions}{The authoritative result and its intended downstream use are identified; unsaved state has not yet been cleared or replaced.}
\contractrow{Inputs}{Recipient and purpose; required molecular format; action/replay needs; calculation, preparation, docking, or campaign evidence to retain.}
\contractrow{Procedure}{Inventory records; export reusable molecular state; export method-specific evidence; create action script if replay is required; verify replay or hashes; describe missing pieces explicitly.}
\contractrow{State mutation}{Export is non-mutating. Script import is destructive to the unsaved scene and must follow export of work that must be retained.}
\contractrow{Evidence}{Exported files, schema and version identifiers, hashes where recorded, replay audit, and a statement of assumptions and unresolved failures.}
\contractrow{Success}{The recipient can identify the starting state, reproduce guarded actions where supported, inspect the method and result, and distinguish accepted decisions from candidates.}
\contractrow{On failure}{Stop on schema, action, binding, hash, or referenced-content mismatch. Do not describe a partial replay as equivalent to a completed audit.}
\contractrow{Boundary}{No autosave, general database, full project bundle, one-click live-session-to-campaign commit, or automatic recovery after Clear, reload, or import.}
\end{skillcontract}

\lead{The appendix in one sentence.}For every task: make the molecular state explicit, choose a method whose scope matches the question, inspect before applying a candidate, retain the evidence needed to reproduce the decision, and stop when the declared scientific boundary is crossed.

\subhead{Capabilities deliberately outside these workflows}

The current workbench does not provide missing-loop construction, membrane or general metal preparation, covalent-ligand parameterization, periodic explicit solvent, PME, production-scale trajectories, binding free energy, global docking, general remote CUDA execution, GFN/xTB, general combinatorial enumeration, or automatic campaign selection. Their absence is part of the operating contract rather than a prompt to approximate them silently.

{\fontsize{8.97}{10.8}\selectfont
This task-centered appendix intentionally omits maintainer file paths, source-line numbers, worker lifecycle details, JSON nesting limits, deployment internals, and the complete repository test list. Those remain valuable implementation and audit documentation, but they are not required for a scientist or agent to understand how to perform and judge the workflows above.\par}

\end{document}
